\section{Microscopic flow, scattering, and kinetic envelopes}
\label{sec:preliminaries}

The microscopic expansion uses globally defined finite-particle flows
and small initial centered correlations. We establish these facts,
identify the collision operator by two-body flux, and derive fixed
spatially summable Gaussian bounds for the assumed kinetic solution
and its spatial gradient. The scattering arguments require neither a
monotone potential nor a regular differential cross-section.

\subsection{Potential assumptions and finite-particle flow}

\begin{assumption}[Potential class]\label{pre:potential}
Write $\Phi(q)=\phi(|q|)$, with $\Phi(q)=0$ for $|q|\geq1$.
Assume either $\Phi\in C^2(\mathbb R^3;\mathbb R)$, or
$\Phi\in C^2(\mathbb R^3\setminus\{0\})$ with
$\Phi(q)\to+\infty$ as $q\to0$.
In the latter case set $\Phi(0)=+\infty$.
Assume thermodynamic stability: for some $B\geq0$,
\begin{equation}\label{eq:pre-stability}
 \sum_{\substack{i,j\in R\\i<j}}\Phi(q_i-q_j)\geq-B|R|
\end{equation}
for every finite configuration of distinct points.
The potential is fixed independently of $\varepsilon$.
\end{assumption}

These are the assumptions of Hypothesis $1'$ in
\cite[Section 8]{PSS14}. Set
\[
 \Phi_\varepsilon(x)=\Phi(x/\varepsilon),\qquad
 \mathscr H_R(Z_R)=\frac12\sum_{i\in R}|v_i|^2+
       \sum_{\substack{i,j\in R\\i<j}}\Phi_\varepsilon(x_i-x_j).
\]
The exact microscopic equations are
\begin{equation}\label{pre:dynamics}
 \dot x_i=v_i,\qquad
 \dot v_i=-\frac1\varepsilon\sum_{j\in R\setminus\{i\}}
             \nabla\Phi((x_i-x_j)/\varepsilon).
\end{equation}

\begin{proposition}[Global finite-particle flow]\label{pre:flow}
For each fixed $\varepsilon>0$ and finite label set $R$,
\eqref{pre:dynamics} has a unique global flow from every finite-energy
initial state in its domain. The flow preserves phase volume and
$\mathscr H_R$. In the singular case, every pair distance along a given
trajectory is bounded away from zero. In the nonsingular case coincident
positions are permitted.
\end{proposition}

\begin{proof}
Taking two particles in \eqref{eq:pre-stability} gives
$\phi(r)\geq-2B$ for $r>0$. In the nonsingular case, approximation by
distinct configurations extends stability to coincident positions.
The vector field is locally Lipschitz on its domain and hence has a
unique maximal local solution. Differentiation gives energy conservation
on its existence interval. For $n=|R|$,
\begin{equation}\label{eq:pre-speed}
 \frac12\sum_{i\in R}|v_i(t)|^2
       \leq\mathscr H_R(Z_R(0))+Bn.
\end{equation}
For $n\geq2$, bounding every pair except one below gives
\begin{equation}\label{eq:pre-pairupper}
 \Phi_\varepsilon(x_i(t)-x_j(t))
 \leq\mathscr H_R(Z_R(0))
             +2B\left(\binom n2-1\right).
\end{equation}
In the singular case this bounds every pair distance away from zero.
Equation \eqref{eq:pre-speed} prevents escape to spatial infinity in
finite time. Thus on each bounded time interval the solution remains in
a compact subset of the vector field's domain and extends beyond any
proposed finite endpoint, forward or backward. In the nonsingular case
one can also use the globally Lipschitz gradient of the compactly
supported $C^2$ potential directly. The cases $n=0,1$ are immediate.

The vector field has zero divergence in $(x_i,v_i)_{i\in R}$:
$\dot x_i$ is independent of $x_i$ and $\dot v_i$ is independent of
$v_i$. Its flow Jacobian therefore solves $\dot J=0$ with $J(0)=1$.
\end{proof}

\subsection{The ensemble and external isolation}

\begin{definition}[Exact initial ensemble]
Let $\mu=\varepsilon^{-2}$ and
\[
 {\cal D}_R^\varepsilon
 =\{Z_R:|x_i-x_j|>\varepsilon\text{ for all }i\ne j\in R\}.
\]
For the prescribed $C^1$ nonnegative probability density $f_0$, set
\begin{align}
 {\cal Z}_\varepsilon
 &=\sum_{n=0}^\infty\frac{\mu^n}{n!}
   \int f_0^{\otimes n}(Z_n)
             \mathbf1_{{\cal D}_{[n]}^\varepsilon}(Z_n)\,dZ_n,
       \label{eq:pre-partition}\\
 W_{0,n}(Z_n)
 &={\cal Z}_\varepsilon^{-1}\mu^n f_0^{\otimes n}(Z_n)
             \mathbf1_{{\cal D}_{[n]}^\varepsilon}(Z_n).
       \label{eq:pre-ensemble}
\end{align}
Here $[n]=\{1,\ldots,n\}$ and empty products equal one.
The probability law on ordered coordinates is
\begin{equation}\label{pre:model}
 \mathbb P(N=n,Z_n\in dZ_n)=\frac1{n!}W_{0,n}(Z_n)\,dZ_n.
\end{equation}
Thereafter the cloud evolves by \eqref{pre:dynamics}.
\end{definition}

Since $1\leq\mathcal Z_\varepsilon\leq e^\mu$, this is a probability law
with finite $N$ almost surely. Initial exclusion gives zero potential
energy and finite kinetic energy. Proposition~\ref{pre:flow} therefore
defines its exact evolution. Let $W_{t,n}$ be the transported density.

For $r=|R|$, the rescaled factorial density is
\begin{equation}\label{pre:factorial}
 F_R^\varepsilon(t,Z_R)=
 \mu^{-r}\sum_{n=0}^\infty\frac1{n!}
       \int W_{t,r+n}(Z_R,Z_{[n]}^{\rm ext})\,dZ_{[n]}^{\rm ext}.
\end{equation}
External labels are distinct from those in $R$.
Symmetry makes the formula independent of their temporary names.
It equals the density defined by the ordered distinct-label sum in the
statement: in a cloud of size $r+n$, there are $(r+n)!/n!$ such selections.

The externally isolated factorial density is
\begin{equation}\label{eq:pre-isolated}
 \widetilde F_R^\varepsilon(t,Z_R)=
 \mu^{-r}\sum_{n=0}^\infty\frac1{n!}
 \int W_{t,r+n}(Z_R,Z_{[n]}^{\rm ext})
 \prod_{\substack{i\in R\\j\in[n]^{\rm ext}}}
       \mathbf1_{\{|x_i-x_j|>\varepsilon\}}\,
             dZ_{[n]}^{\rm ext}.
\end{equation}
Set $F_\varnothing^\varepsilon=\widetilde F_\varnothing^\varepsilon=1$.
External isolation imposes no condition on distances within $R$.
In particular an interacting displayed component can be externally
isolated, and $0\leq\widetilde F_R^\varepsilon\leq F_R^\varepsilon$.
A polymer is a connected component of the instantaneous graph with
edges $|x_i-x_j|\leq\varepsilon$.
The condition in \eqref{eq:pre-isolated} says exactly that the displayed
labels form a union of entire polymers of the full cloud.

\begin{proposition}[Initial insertion and centering]\label{pre:initial}
Put $\rho_0(x)=\int f_0(x,v)\,dv$, and assume $\|\rho_0\|_\infty<\infty$.
At time zero,
\begin{equation}\label{eq:pre-insertion}
 \begin{split}
 \widetilde F_R^\varepsilon(0,Z_R)
 &=F_R^\varepsilon(0,Z_R)\\
 &=f_0^{\otimes R}(Z_R)\mathbf1_{{\cal D}_R^\varepsilon}(Z_R)
       \mathbb E\prod_{i\in R}\mathbf1_{\{N(B_i)=0\}},
 \qquad B_i=\overline B(x_i,\varepsilon).
 \end{split}
\end{equation}
The expectation is over an independent cloud with the original law,
and $N(B)$ counts its positions in $B$.
Consequently $0\leq F_R^\varepsilon(0)\leq f_0^{\otimes R}$.

For $Z_R\in{\cal D}_R^\varepsilon$, define
\begin{equation}\label{eq:pre-initialcenter}
 E_R^\varepsilon(0,Z_R)=
 \sum_{S\subset R}(-1)^{|R\setminus S|}
      f_0^{\otimes(R\setminus S)}(Z_{R\setminus S})
      \widetilde F_S^\varepsilon(0,Z_S).
\end{equation}
There are $c_0>0,C_0<\infty$, independent of $R,\varepsilon$, such that
for sufficiently small $\varepsilon$,
\begin{equation}\label{eq:pre-centersmall}
 |E_R^\varepsilon(0,Z_R)|
       \leq(C_0\varepsilon^{c_0})^{|R|}f_0^{\otimes R}(Z_R).
\end{equation}
One may take $c_0=1/125$.
For fixed $r$,
\begin{equation}\label{eq:pre-mass}
 0\leq1-\int F_R^\varepsilon(0,Z_R)\,dZ_R
       \leq Cr\varepsilon+Cr^2\varepsilon^3,
 \qquad \int F_R^\varepsilon(t)=\int F_R^\varepsilon(0).
\end{equation}
\end{proposition}

\begin{proof}
Insert $r$ particles into \eqref{pre:factorial} at time zero.
Their mutual exclusions give $\mathbf1_{{\cal D}_R^\varepsilon}$.
The remaining partition sum is $\mathcal Z_\varepsilon$ times the
probability that the original cloud misses $\bigcup_iB_i$.
This proves \eqref{eq:pre-insertion}. Initial exclusion already gives
external isolation.

On $\mathcal D_R^\varepsilon$ every subset is separated, so substitution
into \eqref{eq:pre-initialcenter} gives the exact identity
\begin{equation}\label{eq:pre-occupancy}
 E_R^\varepsilon(0,Z_R)
 =(-1)^{|R|}f_0^{\otimes R}(Z_R)
       \mathbb P\{N(B_i)\geq1\text{ for every }i\in R\}.
\end{equation}
Choose a maximal subcollection of centers with mutual distance greater
than $2\varepsilon$. Its balls $B_i$ are disjoint up to null boundaries.
A ball of radius $2\varepsilon$ contains at most $125$ original centers:
their disjoint balls of radius $\varepsilon/2$ fit inside a ball of
radius $5\varepsilon/2$. Hence the chosen subcollection has size
$m\geq |R|/125$. Disjointness ensures that occupants are distinct, and
the factorial upper bound gives
\begin{align*}
 \mathbb P\{N(B_i)\geq1\text{ for every chosen }i\}
 &\leq\mathbb E\prod_{\text{chosen }i}N(B_i)\\
 &=\mu^m\int_{\prod_i(B_i\times\mathbb R^3)}
                          F_m^\varepsilon(0)\,dZ_m\\
 &\leq\bigl(\mu\|\rho_0\|_\infty|B(0,\varepsilon)|\bigr)^m
 \leq(C\varepsilon)^m.
\end{align*}
For $C\varepsilon\leq1$ this proves \eqref{eq:pre-centersmall}.

The loss from mutual exclusion of $r$ independent inserted positions
is at most $Cr^2\varepsilon^3$, by a union bound and bounded $\rho_0$.
For each inserted position,
\[
 \mathbb P\{N(B_i)\geq1\}\leq\mathbb E N(B_i)
 =\mu\int_{B_i\times\mathbb R^3}F_1^\varepsilon(0,z)\,dz
 \leq C\varepsilon.
\]
A second union bound proves the first estimate in \eqref{eq:pre-mass}.
The integral of \eqref{pre:factorial} is
$\mu^{-r}\mathbb E(N)_r$, where
$(N)_r=N(N-1)\cdots(N-r+1)$.
It is finite by the initial factorial bound and conserved since $N$ is
conserved.
\end{proof}

\subsection{Scattering flux and the collision operator}

Use microscopic relative coordinates in this subsection.
For $g\ne0$ and $b\in g^\perp$, $|b|<1$, the incoming free trajectory
meets the interaction sphere at
\[
 n=b-\sqrt{1-|b|^2}\,\frac g{|g|},\qquad n\cdot g<0.
\]
Time translation identifies the asymptotic prescription
$q(\tau)-(b+\tau g)\to0$, $q'(\tau)\to g$ as $\tau\to-\infty$,
with $q(0)=n$, $q'(0)=g$ for $q''=-2\nabla\Phi(q)$.
Let $\tau_*(n,g)\in(0,\infty]$ be its first subsequent exit time from
the open unit ball.
Projection onto the impact disk has Jacobian
\begin{equation}\label{eq:pre-fluxdisk}
 |n\cdot g|\,d\sigma(n)=|g|\,d^2b.
\end{equation}

\begin{lemma}[Residence-time flux]\label{pre:duration}
For every $M<\infty$ there is $C_M<\infty$ such that
\begin{equation}\label{eq:pre-duration}
 \int_{\substack{|g|\leq M\\n\in S^2,\ n\cdot g<0}}
 \tau_*(n,g)|n\cdot g|\,d\sigma(n)\,dg\leq C_M.
\end{equation}
The incoming flux of $\{\tau_*>D,\ |g|\leq M\}$ is at most $C_M/D$.
In particular $\tau_*<\infty$ almost everywhere in incoming flux.
The entrance-to-exit scattering map preserves flux.
\end{lemma}

\begin{proof}
Relative energy is
\[
 \tfrac14|q'|^2+\phi(|q|)=\tfrac14|g|^2.
\]
Since $\phi\geq-2B$, we have $|q'|^2\leq M^2+8B$ throughout the
encounter. Consider
\[
 (n,g,s)\longmapsto(q(s),q'(s)),
 \quad n\cdot g<0,\quad |g|\leq M,\quad0<s<\tau_*(n,g).
\]
Its volume element is $|n\cdot g|\,d\sigma(n)\,dg\,ds$.
At $s=0$, append the velocity vector field to five tangent vectors of
the entrance section. Its normal component in position space is
$n\cdot g$. Phase-volume preservation propagates that determinant.

The map is injective. Two representations of an interior state lie on
the same orbit by uniqueness. A preceding outward crossing cannot be
followed by an inward crossing: outside the ball the path is a
straight line with increasing distance from the origin. If the orbit
repeated an interior state it would be periodic, so it could not have
a past entrance from outside. Hence there is at most one entrance and
elapsed time.

Apply change of variables on compact pieces of the entrance section
and bounded elapsed-time intervals, then exhaust the domain.
The image lies in $B(0,1)\times B(0,\sqrt{M^2+8B})$.
Its finite volume proves \eqref{eq:pre-duration}.
Markov's inequality gives the tail and excludes infinite duration
outside a null set. The same flow-box determinant on the transverse
exit section proves flux preservation. Tangential section data carry
zero flux.
\end{proof}

\begin{definition}[The scattering collision operator]
Write $c=(v+v_*)/2$, $g=v-v_*$, and, for almost every incoming datum,
let $g_{\rm out}$ be the velocity after exit. Set
\[
 v'=c+\tfrac12g_{\rm out},\qquad
 v_*'=c-\tfrac12g_{\rm out}.
\]
For $\psi\in C_c^\infty(\mathbb R^3)$, define $Q_\Phi$ at each $x$ by
\begin{equation}\label{pre:collision}
 \begin{split}
 \int Q_\Phi(f,f)(x,v)\psi(v)\,dv
 =\frac12\int_{\mathbb R^3\times\mathbb R^3}
 & f(x,v)f(x,v_*)|v-v_*|\\
 {}\times\int_{\substack{b\in(v-v_*)^\perp\\|b|<1}}
 &[\,\psi(v')+\psi(v_*')-\psi(v)-\psi(v_*)\,]\,
                         d^2b\,dv\,dv_* .
 \end{split}
\end{equation}
Here $d^2b$ is Euclidean area, and the $g=0$ contribution is zero.
Null exceptional scattering data are ignored.
\end{definition}

\begin{lemma}[Reversibility and normalization]\label{pre:operator}
The measure $d\Lambda=|v-v_*|\,dv\,dv_*\,d^2b$ is preserved by
exchanging incoming and outgoing velocity pairs with the corresponding
entrance parameter. The weak definition \eqref{pre:collision} equals
\begin{equation}\label{eq:pre-Qstrong}
 Q_\Phi(f,f)(x,v)=
 \int_{\mathbb R^3}|v-v_*|
 \int_{\substack{b\in(v-v_*)^\perp\\|b|<1}}
 [\,f(x,v')f(x,v_*')-f(x,v)f(x,v_*)\,]\,d^2b\,dv_* .
\end{equation}
These identities hold whenever the integrals are absolutely convergent,
in particular for Gaussian bounded functions. All scattering branches
are included, and the impact disk has area $\pi$.
\end{lemma}

\begin{proof}
Let $(n_+,g_{\rm out})$ be the exit state of the relative orbit with
entrance $(n,g)$. Since $\Phi$ is even, $-q(\tau_*-s)$ has entrance
$(-n_+,g_{\rm out})$ and exit $(-n,g)$.
Thus $(n,g)\mapsto(-n_+,g_{\rm out})$ is an involution of the incoming
section outside a null set. It preserves flux by
Lemma~\ref{pre:duration} and the reflection $n_+\mapsto-n_+$.
The center velocity is unchanged, and
$(v,v_*)\leftrightarrow(c,g)$ has absolute determinant one.
Equation \eqref{eq:pre-fluxdisk} proves invariance of $d\Lambda$.

Particle exchange also preserves $d\Lambda$, so the two test-function
terms in either pair of \eqref{pre:collision} give equal integrals.
This cancels the factor $1/2$. The scattering involution transfers
the remaining outgoing test to the incoming variable, giving
\eqref{eq:pre-Qstrong}.
Energy conservation gives
$|v'|^2+|v_*'|^2=|v|^2+|v_*|^2$.
The finite disk area and Gaussian integration justify absolute
convergence. No scattering-angle density or choice of one branch is
introduced.
\end{proof}

\subsection{Exceptional deflections and rotational parameters}

\begin{proposition}[Scattering nondegeneracy]\label{pre:scattering}
Outside a set of zero incoming flux, an encounter either is a straight
force-free passage or satisfies
\begin{equation}\label{eq:pre-increments}
 v'\ne v,\qquad v_*'\ne v .
\end{equation}
For $M<\infty$ and $\eta>0$, the non-free relative incoming data $(n,g)$
with $|g|\leq M$
can be restricted to a rotation-invariant compact set, losing at most
$\eta$ relative incoming flux, on which the duration is bounded and
$|g|$, $|v'-v|$, and $|v_*'-v|$ have positive lower bounds.
For fixed radial scattering parameters, each of these nonzero
increments has a uniform spherical directional marginal under common
rotation. The two increments need not be independent.
\end{proposition}

\begin{proof}
Put $G=|g|$, $\rho=|b|$, $J=G\rho$, and $E=G^2/4$.
The cases $G=0$, $J=0$, and $\rho=1$ are null in incoming flux.
For $J>0$, the effective radial potential is
\[
 W_J(r)=\phi(r)+\frac{J^2}{4r^2}.
\]
An incoming turning radius $y$ satisfies $E=W_J(y)$.
A nonsimple turning point also requires $W_J'(y)=0$.
At each fixed $J$, the critical values of the $C^2$ function $W_J$ on
$(0,1)$ are Lebesgue null, by the one-dimensional Sard theorem on a
countable compact exhaustion. The radial incoming flux is
\begin{equation}\label{eq:pre-Sardflux}
 G^3\rho\,dG\,d\rho=2J\,dE\,dJ.
\end{equation}
Fubini excludes nonsimple turning energies up to zero flux.
Lemma~\ref{pre:duration} has already excluded infinite duration.

At a simple incoming turning radius $y\in(0,1)$,
$W_J'(y)<0$ and $E>W_J(r)$ for all $r>y$.
Set $a=J^{-2}$ and
\[
 A_y(r)=y^{-2}-r^{-2},\qquad
 C_y(r)=4(\phi(y)-\phi(r)),\qquad
 H_{y,a}(r)=A_y(r)+aC_y(r).
\]
The half angular sweep is
\begin{equation}\label{eq:pre-halfangle}
 {\cal A}_y(a)=\int_y^\infty\frac{dr}{r^2\sqrt{H_{y,a}(r)}}.
\end{equation}
Indeed $\dot r^2=J^2H_{y,a}(r)$ and the angular speed is $J/r^2$.
For admissible $a$, one has $H_{y,a}(r)>0$ for $r>y$,
$\partial_rH_{y,a}(y)>0$, and
$\lim_{r\to\infty}H_{y,a}(r)>0$.
The ratios $C_y/A_y$ and $H_{y,a}/A_y$ extend continuously to
$[y,\infty]$. At $y$ use the ratios of first derivatives; at infinity
use finite range. The second ratio is strictly positive on this compact
interval, so it has a positive minimum, while $C_y/H_{y,a}$ is bounded.
These properties also show that the admissible set of $a$ is open.

For $a'$ near $a$, the binomial expansion in
\[
 H_{y,a'}^{-1/2}
 =H_{y,a}^{-1/2}
       \left(1+(a'-a)\frac{C_y}{H_{y,a}}\right)^{-1/2}
\]
converges uniformly after multiplication by the integrable majorant
$r^{-2}H_{y,a}^{-1/2}$. Thus $\mathcal A_y$ is real analytic in the
parameter $a$, although $\phi$ is only $C^2$.
Twice differentiating gives
\begin{equation}\label{eq:pre-convex}
 {\cal A}_y''(a)=\frac34
       \int_y^\infty\frac{C_y(r)^2\,dr}{r^2H_{y,a}(r)^{5/2}}.
\end{equation}
Near $y$ the integrand is $O((r-y)^{-1/2})$, since $C_y(r)=O(r-y)$
and $H_{y,a}$ has a simple root. It is also integrable at infinity.
It is strictly positive unless $\phi(r)=\phi(y)$ for every $r\geq y$.
Finite range makes this constant zero; the encountered orbit is then
a straight force-free passage.

For a non-free encounter $\mathcal A_y$ is strictly convex.
The deflection angle is $\pi-2\mathcal A_y$ modulo $2\pi$.
Therefore $g_{\rm out}=g$ or $g_{\rm out}=-g$ requires
$\mathcal A_y\in(\pi/2)\mathbb Z$.
Each such level has at most two preimages on an admissible interval.
There are countably many levels and countably many open intervals.
On a simple turning branch, $(y,J)\mapsto(E,J)$, $E=W_J(y)$,
is a local $C^1$ diffeomorphism.
A countable cover by these charts, Fubini in $(y,a)$, and
\eqref{eq:pre-Sardflux} prove flux nullness.
These exceptional equalities are exactly the two failures of
\eqref{eq:pre-increments}.

For the compact-set assertion, exhaust the simple-scattering non-free
radial parameter domain by compact subsets avoiding $G=0$, $J=0$,
$\rho=1$, and the zero-increment levels.
This covers all non-free data outside a null set.
On each such compact subset, continuous dependence for the $C^1$
ODE flow and transversality of the finite exit give a continuous
exit time and scattering map. Thus durations are bounded and all the
listed nonzero quantities have positive minima.
The relative flux measure on $G\leq M$ is finite, so an exhaustion
member loses at most $\eta$ flux. Lifting with all rotations makes
the set rotation-invariant.

Finally, write $(g,b)=(GRe_1,\rho Re_2)$, $R\in SO(3)$.
At fixed $(G,\rho)$ the orientation factor in incoming flux is a
constant multiple of Haar measure on $SO(3)$.
Central-force rotation equivariance maps each nonzero increment to
$R$ times a fixed vector. Its direction has uniform measure on $S^2$.
This is a marginal statement and leaves the joint pair correlated.
\end{proof}

\subsection{A fixed spatial Gaussian envelope on the kinetic interval}

Use the norm $\|\cdot\|_{\mathrm{Bol},\alpha}$ defined in
\eqref{intro:norm}, and put
\[
 \|h\|_{\infty,\alpha}
       =\sup_{x,v}e^{\alpha|v|^2}|h(x,v)|.
\]
For vector-valued functions use Euclidean norm.
The following proof adapts the weighted transport estimate in
\cite[Proposition A.1]{DHM25}. Only finite impact area and kinetic
energy conservation are needed, so the argument applies to $Q_\Phi$.

\begin{proposition}[Known-solution envelopes]\label{pre:envelopes}
Suppose $f$ is the classical nonnegative solution in the statement, and
\[
 A=\sup_{0\leq t\leq T}\|f(t)\|_{\infty,2\beta}<\infty,\qquad
 B_0=\|f_0\|_{\mathrm{Bol},2\beta}
               +\|\nabla_xf_0\|_{\mathrm{Bol},2\beta}<\infty.
\]
Then
\begin{equation}\label{eq:pre-envelope}
 \sup_{0\leq t\leq T}
 \left(\|f(t)\|_{\mathrm{Bol},\beta}
       +\|\nabla_xf(t)\|_{\mathrm{Bol},\beta}\right)
       \leq C(A,B_0,\beta,T)<\infty.
\end{equation}
Fix, once and for all,
\begin{equation}\label{eq:pre-beta0}
 \beta_0=\beta/2.
\end{equation}
Thus all independent inputs $f(t)$ and their spatial Lipschitz
envelopes have a common bound with the fixed margin
$\beta-\beta_0$. Constants in \eqref{eq:pre-envelope} are independent
of the subsequent time-layer subdivision.
\end{proposition}

\begin{proof}
Set
\[
 \gamma(t)=\frac{3\beta}{2}-\frac{\beta t}{2T},\qquad
 c=\frac{\beta}{2T},\qquad M(v)=e^{-2\beta|v|^2},
\]
and $(S(r)h)(x,v)=h(x-rv,v)$.
Notice $\beta\leq\gamma(t)\leq3\beta/2$.
Define a positive linear operator for $h\geq0$ by
\begin{equation}\label{eq:pre-positiveB}
 \begin{split}
 ({\cal B}h)(x,v)=\int_{\mathbb R^3}|v-v_*|\int_{|b|<1}
 \bigl[&M(v_*')h(x,v')+M(v')h(x,v_*')\\
       &+M(v_*)h(x,v)+M(v)h(x,v_*)\bigr]\,d^2b\,dv_* .
 \end{split}
\end{equation}
The disk is in $(v-v_*)^\perp$.

We first prove that for $0\leq s<t\leq T$, $r=t-s$,
\begin{align}
 \|S(r){\cal B}h\|_{\mathrm{Bol},\gamma(t)}
       &\leq Cr^{-1/2}\|h\|_{\mathrm{Bol},\gamma(s)},
                    \label{eq:pre-VolterraBol}\\
 \|S(r){\cal B}h\|_{\infty,\gamma(t)}
       &\leq Cr^{-1/2}\|h\|_{\infty,\gamma(s)}.
                    \label{eq:pre-Volterrainfty}
\end{align}
Let $H_s(y)=\sup_u e^{\gamma(s)|u|^2}h(y,u)$.
For each product in \eqref{eq:pre-positiveB}, energy conservation,
or its identity version for a loss product, gives
\[
 M(u_*)h(y,u)
       \leq H_s(y)e^{-\gamma(s)(|v|^2+|v_*|^2)}.
\]
After multiplication by $e^{\gamma(t)|v|^2}$ this is at most
\[
 H_s(y)e^{-\beta|v_*|^2}
                    e^{-cr(|v|^2+|v_*|^2)}.
\]
Since
\[
 |v-v_*|e^{-cr(|v|^2+|v_*|^2)/2}\leq C_c r^{-1/2},
\]
the disk area $\pi$ and the Gaussian integral in $v_*$ imply
\begin{equation}\label{eq:pre-spatialmajorant}
 e^{\gamma(t)|v|^2}(S(r){\cal B}h)(x,v)
       \leq Cr^{-1/2}e^{-cr|v|^2/2}H_s(x-rv).
\end{equation}
This proves \eqref{eq:pre-Volterrainfty}.

For the summed spatial estimate put
$a_k=\sup_{|y-k|\leq1}H_s(y)$.
If $m\leq r|v|<m+1$ and $|x-k|\leq1$, then
$|x-rv-k|\leq m+2$, whence
\[
 H_s(x-rv)\leq
       \sum_{\substack{\ell\in\mathbb Z^3\\|\ell-k|\leq m+3}}a_\ell.
\]
Indeed lattice-centered unit balls cover $\mathbb R^3$.
For $r\leq T$, the damping factor on that shell is at most
$e^{-cm^2/(2T)}$.
Bound the supremum over $v$ by summation over $m\geq0$ and then sum
over $k$. Each $a_\ell$ occurs at most $C(m+4)^3$ times.
Convergence of
$\sum_{m\geq0}e^{-cm^2/(2T)}(m+4)^3$ proves
\eqref{eq:pre-VolterraBol}.
The same shell argument, now using
\[
 e^{\gamma(t)|v|^2}|S(t)h_0(x,v)|
 \leq e^{-ct|v|^2}
       \sup_u e^{\gamma(0)|u|^2}|h_0(x-tv,u)|,
\]
gives
\begin{equation}\label{eq:pre-freebound}
 \sup_{0\leq t\leq T}\|S(t)h_0\|_{\mathrm{Bol},\gamma(t)}
       \leq C\|h_0\|_{\mathrm{Bol},\gamma(0)}.
\end{equation}
At $t=0$ this is immediate.

Define
\[
 ({\cal V}h)(t)=A\int_0^t S(t-s){\cal B}h(s)\,ds.
\]
The characteristic equation, $f\geq0$, and $f\leq AM$ imply
\begin{equation}\label{eq:pre-linearcomparison}
 0\leq f(t)\leq S(t)f_0+({\cal V}f)(t).
\end{equation}
Its $n$-fold iteration is the finite inequality
\[
 f\leq\sum_{j=0}^{n-1}{\cal V}^j(S(\cdot)f_0)+{\cal V}^n f.
\]
No finiteness of the unknown summed spatial norm has been assumed.
Using \eqref{eq:pre-Volterrainfty} and
$\|f(t)\|_{\infty,\gamma(t)}\leq A$ yields
\begin{equation}\label{eq:pre-remaindergamma}
 \|({\cal V}^n f)(t)\|_{\infty,\gamma(t)}
       \leq A\frac{(CA\Gamma(1/2))^nt^{n/2}}
                           {\Gamma(1+n/2)}.
\end{equation}
This tends to zero uniformly for $t\leq T$.
The time-convolution formula follows inductively from
\[
 \int_0^t(t-s)^{-1/2}s^{n/2}\,ds
 =t^{(n+1)/2}
       \frac{\Gamma(1/2)\Gamma(1+n/2)}
            {\Gamma(1+(n+1)/2)}.
\]
Equations \eqref{eq:pre-VolterraBol} and
\eqref{eq:pre-freebound} bound the summed spatial norm of the
$j$th series term by
\[
 C B_0\frac{(CA\Gamma(1/2))^jT^{j/2}}{\Gamma(1+j/2)}.
\]
The sum is finite. Letting $n\to\infty$ in the pointwise comparison
and then taking the summed spatial norm proves the required bound
for $f$ at exponent $\gamma(t)$.

To handle the derivative without assuming its envelope, fix a
coordinate unit vector $e$ and $0<|\delta|\leq1$, and put
\[
 d_\delta(t,x,v)
       =\frac{f(t,x+\delta e,v)-f(t,x,v)}{\delta}.
\]
Use $a'b'-ab=(a'-a)b'+a(b'-b)$ for both gain and loss products in
\eqref{eq:pre-Qstrong}. Each undifferenced factor is a spatial
translate of $f$, bounded by $AM$, so
\[
 |d_\delta(t)|
       \leq S(t)|d_\delta(0)|+{\cal V}|d_\delta|(t).
\]
For fixed $\delta$,
$\|d_\delta(t)\|_{\infty,\gamma(t)}\leq2A/|\delta|$.
The remainder argument \eqref{eq:pre-remaindergamma} therefore
applies and vanishes as $n\to\infty$, before taking $\delta\to0$.
The fundamental theorem of calculus gives
\[
 |d_\delta(0,x,v)|
       \leq\int_0^1|\partial_e f_0(x+\theta\delta e,v)|\,d\theta.
\]
Translations of length at most one change the summed spatial
supremum norm by at most an absolute factor, since each translated
unit ball is covered by a fixed number of lattice-centered unit
balls and each covering ball occurs a bounded number of times.
Hence
$\|d_\delta(0)\|_{\mathrm{Bol},\gamma(0)}\leq CB_0$,
uniformly in $\delta$.
The same series estimate bounds
$\|d_\delta(t)\|_{\mathrm{Bol},\gamma(t)}$ uniformly in $t,\delta$.
Pointwise convergence to $\partial_e f$, followed by Fatou's
inequality for the sum of local suprema, gives the derivative
bound. Sum over the three coordinate directions and use
$\gamma(t)\geq\beta$ to obtain \eqref{eq:pre-envelope}.
\end{proof}

\begin{remark}\label{pre:no-manybody-conclusion}
Lemma~\ref{pre:duration} concerns isolated two-body encounters, not
dynamically formed many-particle bound components.
Proposition~\ref{pre:scattering} supplies rotational marginals, not
independence from a trajectory that already depends on the same
parameters. The history estimates must establish the required
many-particle and conditional parameter statements separately.
\end{remark}
